\section{Traversal of quasi-subsequential finite-state transducers}

\Cref{alg:qsfst_traversal} represents a deterministic procedure that translates an
input word using a $k$-separated quasi-subsequential finite-state transducer. The
output of the procedure is one of the possible translations of the input word.

\begin{algorithm}[h]
  \DontPrintSemicolon
  \KwIn{$\langle \Sigma, \Omega, Q, X, s, F, \tau, \rho, \eta, \eta_x, \lambda, \psi
      \rangle, \sigma_1 \sigma_2 \ldots \sigma_n$}

  $o \gets \epsilon$\;\nllabel{alg:qsfst:init_begin}
  $q \gets s$\;\nllabel{alg:qsfst:init_end}
  \For {$i \in \{ 1, 2, \ldots, n \}$} {
    $L \gets \langle \rangle$\;\nllabel{alg:qsfst:l_init}
    \For {$j \in \{ 1, 2, \ldots, k \}$} {
      $t_j \gets \tau_j(q_j, \sigma_i)$\;\nllabel{alg:qsfst:t}
      $r_j \gets \rho_j(t_j)$\;\nllabel{alg:qsfst:r}
      \If {$\lambda_j(q_j, \sigma_i)$ is \textbf{defined}} {
        $L \gets L \cdot \langle \lambda_j(q_j, \sigma_i) \rangle$\;\nllabel{alg:qsfst:l}
      }
    }
    \If {$L = \langle \rangle$} {
      \Return \textbf{undefined}\;\nllabel{alg:qsfst:undefined}
    }
    $o \gets o \cdot lcp(L)$\;\nllabel{alg:qsfst:lcp}
    $x \gets r_1$\;\nllabel{alg:qsfst:x_begin}
    \For {$j \in \{ 2, 3, \ldots, k \}$} {
      $x \gets \eta_x(x, r_j)$\;\nllabel{alg:qsfst:x_end}
    }
    \For {$j \in \{ 1, 2, \ldots, k \}$} {
      $q_j \gets \eta_j(t_j, x)$\;\nllabel{alg:qsfst:eta}
    }
  }
  \For {$i \in \{ 1, 2, \ldots, k \}$} {
    \If {$q_i \in F_i$} {
      \Return $o \cdot \psi_i(q_i)$\;\nllabel{alg:qsfst:final}
    }
  }
  \Return \textbf{undefined}\;

  \caption{Traversal of a $k$-separated quasi-subsequential finite-state transducer.}
  \label{alg:qsfst_traversal}
\end{algorithm}

Here, $L$ is the list of the defined outputs $\lambda_j(q_j, \sigma_i)$, and
$lcp(L)$ is their longest common prefix, computed by \Cref{alg:lcp} below. By
definition, at the beginning of the $i$-th iteration of the outer loop $q =
\delta^*(s, \sigma_1 \ldots \sigma_{i-1})$ and $o = \lambda^*(s, \sigma_1 \ldots
\sigma_{i-1})$, and line~\ref{alg:qsfst:undefined} is reached if and only if
$\lambda(q, \sigma_i)$ is undefined. Therefore, if \Cref{alg:qsfst_traversal} returns
$v$, then $\langle \sigma_1 \sigma_2 \ldots \sigma_n, v \rangle \in
\mathcal{R}_\mathcal{T}$, and if $\sigma_1 \sigma_2 \ldots \sigma_n$ has a translation,
then \Cref{alg:qsfst_traversal} returns one.

The complexity bound for \Cref{alg:qsfst_traversal} depends on the ability to compute
longest common prefixes of at most $k$ words in time proportional to $k$ times the
length of the result. \Cref{alg:lcp} represents one possible way of doing that: it
compares the words column by column and stops at the first mismatch.

\begin{algorithm}[h]
  \DontPrintSemicolon
  \KwIn{$w_1, w_2, \ldots, w_k$}

  $n \gets |w_1|$\;\nllabel{alg:lcp:min_begin}
  \For {$j \in \{ 2, 3, \ldots, k \}$} {
    \If {$|w_j| < n$} {
      $n \gets |w_j|$\;\nllabel{alg:lcp:min_end}
    }
  }

  $p \gets \epsilon$\;
  \For {$i \in \{ 1, 2, \ldots, n \}$} {
    \For {$j \in \{ 2, 3, \ldots, k \}$} {
      \If {$w_j[i] \neq w_{j-1}[i]$} {\nllabel{alg:lcp:compare}
        \Return $p$
      }
    }

    $p \gets p \cdot w_1[i]$\;\nllabel{alg:lcp:append}
  }

  \Return $p$

  \caption{Longest common prefix of $k$ words.}
  \label{alg:lcp}
\end{algorithm}

\begin{proposition}\label{prop:lcp_complexity}
  Let $w_1, w_2, \ldots, w_k$ be words over an alphabet $\Sigma$ for $k > 0$ and let
  $n = \min_{i=1}^k |w_i|$. Then \Cref{alg:lcp} returns $p = \bigwedge_{i=1}^k w_i$
  on input $w_1, w_2, \ldots, w_k$ and completes in $\mathcal{O}(k(|p| + 1))$ time.
  In particular, it completes in $\mathcal{O}(k(n + 1))$ time.
\end{proposition}

\begin{proof}
  At the beginning of the $i$-th iteration of the outer loop, $p = w_1[1] \ldots
  w_1[i-1]$ is a common prefix of $w_1, w_2, \ldots, w_k$. The iteration either
  returns $p$ because two of the words differ at position $i$, in which case $p$ is
  the longest common prefix, or extends $p$ by the common symbol at position $i$. If
  the loop completes, then $p = w_1[1] \ldots w_1[n]$ is a common prefix of length
  $n$, which is the longest possible. Therefore, $p = \bigwedge_{i=1}^k w_i$.
  Lines~\ref{alg:lcp:min_begin}--\ref{alg:lcp:min_end} require $\mathcal{O}(k)$ time.
  The outer loop completes at most $|p|$ iterations and performs at most one further,
  interrupted iteration. Each iteration executes line~\ref{alg:lcp:compare} at most
  $k - 1$ times and line~\ref{alg:lcp:append} at most once, each in $\mathcal{O}(1)$
  time, assuming that concatenation with a single symbol requires $\mathcal{O}(1)$
  time. Therefore, the total runtime complexity bound for \Cref{alg:lcp} is
  $\mathcal{O}(k(|p| + 1))$. Since $p$ is a prefix of every $w_i$, $|p| \leq n$.
\end{proof}

The term $\mathcal{O}(k)$ cannot be omitted: lines~\ref{alg:lcp:min_begin}--%
\ref{alg:lcp:min_end} require $k - 1$ steps even if $n = 0$.

\begin{proposition}
  Let $\mathcal{T} = \langle \Sigma, \Omega, Q, X, s, F, \tau, \rho, \eta, \eta_x,
  \lambda, \psi \rangle$ be a $k$-separa\-ted quasi-subsequential finite-state
  transducer. Let $u = \sigma_1 \sigma_2 \ldots \sigma_n$ be a word of length $n \in
  \mathbb{N}$ over $\Sigma$ such that $\langle u, v \rangle \in \mathcal{R}_\mathcal{T}$
  for some $v \in \Omega^*$, and let $m = \max \{ |v| \mid \exists v \in \Omega^* (
  \langle u, v \rangle \in \mathcal{R}_\mathcal{T}) \}$. Then
  \Cref{alg:qsfst_traversal} completes on input $\mathcal{T}$ and $u$ in
  $\mathcal{O}(k(n+m))$ time.
\end{proposition}

\begin{proof}
  Let $p_i$ be the longest common prefix computed at line~\ref{alg:qsfst:lcp} in the
  $i$-th iteration of the outer loop. Since $u$ has a translation,
  line~\ref{alg:qsfst:undefined} is never reached, and \Cref{alg:qsfst_traversal}
  returns $v = p_1 p_2 \ldots p_n \psi_j(q_j)$ for some $j \in \{ 1, 2, \ldots, k \}$.
  Thus, $\sum_{i=1}^n |p_i| \leq |v| \leq m$.
  Lines~\ref{alg:qsfst:init_begin}--\ref{alg:qsfst:init_end} are executed once and
  require $\mathcal{O}(1)$ time.
  Line~\ref{alg:qsfst:l_init} is executed $n$ times and requires $\mathcal{O}(1)$ time.
  Lines~\ref{alg:qsfst:t}--\ref{alg:qsfst:l} are executed $kn$ times and require
  $\mathcal{O}(1)$ time.
  Line~\ref{alg:qsfst:lcp} is executed $n$ times. By \Cref{prop:lcp_complexity}, its
  $i$-th execution requires $\mathcal{O}(k(|p_i| + 1))$ time, assuming that
  concatenation of words requires time proportional to the length of the appended
  word. Over all iterations, this is $\mathcal{O}(k(n + \sum_{i=1}^n |p_i|)) \subseteq
  \mathcal{O}(k(n + m))$ time.
  Lines~\ref{alg:qsfst:x_begin}--\ref{alg:qsfst:x_end} are executed at most $kn$
  times and require $\mathcal{O}(1)$ time.
  Line~\ref{alg:qsfst:eta} is executed $kn$ times and requires $\mathcal{O}(1)$ time.
  The final loop is executed at most $k$ times, and line~\ref{alg:qsfst:final}
  requires $\mathcal{O}(m)$ time.
  Therefore, the total runtime complexity bound for \Cref{alg:qsfst_traversal} is
  $\mathcal{O}(k(n+m))$.
\end{proof}

The early termination of \Cref{alg:lcp} is essential for this bound. A bound of
$\mathcal{O}(k \min_j |\lambda_j(q_j, \sigma_i)|)$ per execution of
line~\ref{alg:qsfst:lcp} would not suffice, because the shortest output of a
component can be much longer than the common prefix that is appended to $o$. If $u$
has no translation, the same argument bounds the running time by $\mathcal{O}(k(n +
|o|))$, where $o$ is the output computed before the procedure returns.
